%% kapa: exam=1 round=20 year=2009 period=2 form=A subject=accounting subject_ko=회계학 seq=9 number=9
%% source: Q-Net 감정평가사 기출문제 (cst00309) artlSeq=5024178 file=extracted/1008/5024178/2009년 제20회 감정평가사 1차시험 기출문제/2교시.hwp
%% conversion: kordoc 4.24.51 -> tools/gen1.py (auto); review=proofread-sonnet; flags=none
\begin{question}{9}
다음 중 유형자산의 취득원가에 대한 회계처리의 내용으로 옳지 \uline{않은} 것은?
\begin{choices}
  \item 자산의 취득, 건설, 개발에 따른 복구비용이 자산의 내용연수에 걸쳐 발생하는 경우에는 해당 기간의 비용 또는 제조원가로 처리한다.
  \item 유형자산을 취득 또는 사용가능한 상태로 준비하는 과정과 직접 관련이 없는 일반관리비나 경비 등은 유형자산의 취득원가에 포함하지 않는다.
  \item 유형자산을 취득하는 과정에서 매입할인이 있는 경우에는 이를 차감하여 취득원가를 산출한다.
  \item 유형자산을 장기후불조건으로 구입하거나, 대금지급기간이 일반적인 신용기간보다 긴 경우 취득원가는 취득시점의 현금구입가격으로 한다.
  \item 자가건설에 따른 내부이익과 자가건설 과정에서 원재료, 인력 및 기타 자원의 낭비로 인한 비정상적인 원가도 자가건설한 유형자산의 취득원가에 포함한다.
\end{choices}
\end{question}
\answer{5}
